\section{The cycle decomposition}

\begin{theorem}
[disjoint cycle decomposition of permutations]
\label{thm.perm.dcd.main}
\lean{AlgebraicCombinatorics.CycleDecomposition.exists_dcd, AlgebraicCombinatorics.CycleDecomposition.dcd_unique_cycleType, AlgebraicCombinatorics.CycleDecomposition.canonicalDcd_exists_unique}
\leantarget
\leanok
Let $X$
be a finite set. Let $\sigma$ be a permutation of $X$. Then: \medskip

\textbf{(a)} There is a list%
\begin{align*}
&  \Big(\left(  a_{1,1},a_{1,2},\ldots,a_{1,n_{1}}\right)  ,\\
&  \ \ \ \left(  a_{2,1},a_{2,2},\ldots,a_{2,n_{2}}\right)  ,\\
&  \ \ \ \ldots,\\
&  \ \ \ \left(  a_{k,1},a_{k,2},\ldots,a_{k,n_{k}}\right)  \Big)
\end{align*}
of nonempty lists of elements of $X$ such that:

\begin{itemize}
\item each element of $X$ appears exactly once in the composite list%
\begin{align*}
&  (a_{1,1},a_{1,2},\ldots,a_{1,n_{1}},\\
&  \ \ \ a_{2,1},a_{2,2},\ldots,a_{2,n_{2}},\\
&  \ \ \ \ldots,\\
&  \ \ \ a_{k,1},a_{k,2},\ldots,a_{k,n_{k}}),
\end{align*}
and

\item we have%
\[
\sigma=\operatorname*{cyc}\nolimits_{a_{1,1},a_{1,2},\ldots,a_{1,n_{1}}}%
\circ\operatorname*{cyc}\nolimits_{a_{2,1},a_{2,2},\ldots,a_{2,n_{2}}}%
\circ\cdots\circ\operatorname*{cyc}\nolimits_{a_{k,1},a_{k,2},\ldots
,a_{k,n_{k}}}.
\]

\end{itemize}

Such a list is called a \emph{disjoint cycle decomposition} (or short
\emph{DCD}) of $\sigma$. Its entries (which themselves are lists of elements
of $X$) are called the \emph{cycles} of $\sigma$. \medskip

\textbf{(b)} Any two DCDs of $\sigma$ can be obtained from each other by
(repeatedly) swapping the cycles with each other, and rotating each cycle
(i.e., replacing $\left(  a_{i,1},a_{i,2},\ldots,a_{i,n_{i}}\right)  $ by
$\left(  a_{i,2},a_{i,3},\ldots,a_{i,n_{i}},a_{i,1}\right)  $). \medskip

\textbf{(c)} Now assume that $X$ is a set of integers (or, more generally, any
totally ordered finite set). Then, there is a unique DCD
\begin{align*}
&  \Big(\left(  a_{1,1},a_{1,2},\ldots,a_{1,n_{1}}\right)  ,\\
&  \ \ \ \left(  a_{2,1},a_{2,2},\ldots,a_{2,n_{2}}\right)  ,\\
&  \ \ \ \ldots,\\
&  \ \ \ \left(  a_{k,1},a_{k,2},\ldots,a_{k,n_{k}}\right)  \Big)
\end{align*}
of $\sigma$ that satisfies the additional requirements that

\begin{itemize}
\item we have $a_{i,1}\leq a_{i,p}$ for each $i\in\left[  k\right]  $ and each
$p\in\left[  n_{i}\right]  $ (that is, each cycle in this DCD is written with
its smallest entry first), and

\item we have $a_{1,1}>a_{2,1}>\cdots>a_{k,1}$ (that is, the cycles appear in
this DCD in the order of decreasing first entries).
\end{itemize}
\end{theorem}

\begin{proof}
\leanok
\textbf{(a)} Let $\mathcal{D}$ be the cycle digraph of $\sigma$.
This cycle digraph $\mathcal{D}$ has the following two properties:

\begin{itemize}
\item \textit{Outbound uniqueness:} For each node $i$, there is exactly one
arc outgoing from $i$. (Indeed, this is the arc from $i$ to $\sigma\left(
i\right)  $.)

\item \textit{Inbound uniqueness:} For each node $i$, there is exactly one arc
incoming into $i$. (Indeed, this is the arc from $\sigma^{-1}\left(  i\right)
$ to $i$, since $\sigma$ is a permutation and therefore invertible.)
\end{itemize}

Using these two properties, we will now show that the cycle digraph
$\mathcal{D}$ consists of several node-disjoint cycles.

Indeed, if two cycles $C$ and $D$ of
$\mathcal{D}$ have a node in common, then they are identical
(because the outbound
uniqueness property prevents these cycles from ever separating after meeting
at the common node). In other words, any two cycles $C$ and $D$ of
$\mathcal{D}$ are either identical or node-disjoint.

Now, let $i$ be any node of $\mathcal{D}$. Then, if we start at $i$ and follow
the outgoing arcs, then we obtain an infinite walk%
\[
\sigma^{0}\left(  i\right)  \rightarrow\sigma^{1}\left(  i\right)
\rightarrow\sigma^{2}\left(  i\right)  \rightarrow\sigma^{3}\left(  i\right)
\rightarrow\cdots
\]
along our digraph $\mathcal{D}$. Since $X$ is finite, the Pigeonhole Principle
guarantees that this walk will eventually revisit a node it has already been
to; i.e., there exist two integers $u,v\in\mathbb{N}$ with $u<v$ and
$\sigma^{u}\left(  i\right)  =\sigma^{v}\left(  i\right)  $. Let us pick two
such integers $u,v$ with the \textbf{smallest} possible $v$. Thus,
\begin{equation}
\text{the }v\text{ nodes }\sigma^{0}\left(  i\right)  ,\sigma^{1}\left(
i\right)  ,\ldots,\sigma^{v-1}\left(  i\right)  \text{ are distinct}
\label{pf.thm.perm.dcd.main.a.dist}%
\end{equation}
(since otherwise, $v$ would not be smallest possible). Now, $\sigma$ is a
permutation and thus has an inverse $\sigma^{-1}$. Applying the map
$\sigma^{-1}$ to both sides of the equality $\sigma^{u}\left(  i\right)
=\sigma^{v}\left(  i\right)  $, we obtain $\sigma^{u-1}\left(  i\right)
=\sigma^{v-1}\left(  i\right)  $. However, if we had $u\geq1$, then
(\ref{pf.thm.perm.dcd.main.a.dist}) would entail $\sigma^{u-1}\left(
i\right)  \neq\sigma^{v-1}\left(  i\right)  $,
which would contradict $\sigma^{u-1}\left(  i\right)
=\sigma^{v-1}\left(  i\right)  $. Thus, we cannot have $u\geq1$. Hence, $u<1$,
so that $u=0$. Therefore, $\sigma^{u}\left(  i\right)  =\sigma^{0}\left(
i\right)  $, so that $\sigma^{0}\left(  i\right)  =\sigma^{u}\left(  i\right)
=\sigma^{v}\left(  i\right)  $. This shows that our walk is
circular: it comes back to its starting node $\sigma^{0}\left(  i\right)  =i$
after $v$ steps. We have thus found a cycle in our digraph $\mathcal{D}$:%
\[
\sigma^{0}\left(  i\right)  \rightarrow\sigma^{1}\left(  i\right)
\rightarrow\sigma^{2}\left(  i\right)  \rightarrow\cdots\rightarrow\sigma
^{v}\left(  i\right)  =\sigma^{0}\left(  i\right)  .
\]
This shows that the node $i$ lies on a
cycle $C_{i}$ of $\mathcal{D}$.

We thus have shown that each node $i$ of
$\mathcal{D}$ lies on a cycle $C_{i}$.
Any arc of our digraph $\mathcal{D}$ must belong to one of these chosen
cycles.

Combining these facts, we conclude that $\mathcal{D}$ consists of several
node-disjoint cycles. Let us label these cycles as%
\begin{align*}
&  \left(  a_{1,1}\rightarrow a_{1,2}\rightarrow\cdots\rightarrow a_{1,n_{1}%
}\rightarrow a_{1,1}\right)  ,\\
&  \left(  a_{2,1}\rightarrow a_{2,2}\rightarrow\cdots\rightarrow a_{2,n_{2}%
}\rightarrow a_{2,1}\right)  ,\\
&  \ldots,\\
&  \left(  a_{k,1}\rightarrow a_{k,2}\rightarrow\cdots\rightarrow a_{k,n_{k}%
}\rightarrow a_{k,1}\right)
\end{align*}
(making sure to label each cycle only once). Then, each element of $X$ appears
exactly once in the composite list%
\begin{align*}
&  (a_{1,1},a_{1,2},\ldots,a_{1,n_{1}},\\
&  \ \ \ a_{2,1},a_{2,2},\ldots,a_{2,n_{2}},\\
&  \ \ \ \ldots,\\
&  \ \ \ a_{k,1},a_{k,2},\ldots,a_{k,n_{k}}),
\end{align*}
and we have%
\[
\sigma=\operatorname*{cyc}\nolimits_{a_{1,1},a_{1,2},\ldots,a_{1,n_{1}}}%
\circ\operatorname*{cyc}\nolimits_{a_{2,1},a_{2,2},\ldots,a_{2,n_{2}}}%
\circ\cdots\circ\operatorname*{cyc}\nolimits_{a_{k,1},a_{k,2},\ldots
,a_{k,n_{k}}}%
\]
(since $\sigma$ moves any node $i\in X$ one step forward along its chosen
cycle). This proves Theorem \ref{thm.perm.dcd.main} \textbf{(a)}.

\medskip

\textbf{(b)} See \cite[Theorem 1.5.3]{Goodman} or \cite[Chapter I, \S 5.7,
Proposition 7]{Bourba74}. Let%
\begin{align*}
&  \Big(\left(  a_{1,1},a_{1,2},\ldots,a_{1,n_{1}}\right)  ,\\
&  \ \ \ \left(  a_{2,1},a_{2,2},\ldots,a_{2,n_{2}}\right)  ,\\
&  \ \ \ \ldots,\\
&  \ \ \ \left(  a_{k,1},a_{k,2},\ldots,a_{k,n_{k}}\right)  \Big)
\end{align*}
be a DCD of $\sigma$. Then, for each $i\in X$, the cycle of this DCD that
contains $i$ is uniquely determined by $\sigma$ and $i$ up to cyclic rotation
(indeed, it is a rotated version of the list $\left(  i,\sigma\left(
i\right)  ,\sigma^{2}\left(  i\right)  ,\ldots,\sigma^{r-1}\left(  i\right)
\right)  $, where $r$ is the smallest positive integer satisfying $\sigma
^{r}\left(  i\right)  =i$). Therefore, all cycles of this DCD are uniquely
determined by $\sigma$ up to cyclic rotation and up to the relative order in
which these cycles appear in the DCD. This is the claim of
Theorem \ref{thm.perm.dcd.main} \textbf{(b)}.

\medskip

\textbf{(c)} In order to obtain a DCD of $\sigma$ that satisfies these two
requirements, it suffices to

\begin{itemize}
\item start with an arbitrary DCD of $\sigma$,

\item then rotate each cycle of this DCD so that it begins with its smallest
entry, and

\item then repeatedly swap these cycles so they appear in the order of
decreasing first entries.
\end{itemize}

It is clear that the result of this procedure is uniquely determined (a
consequence of Theorem \ref{thm.perm.dcd.main} \textbf{(b)}). Thus, Theorem
\ref{thm.perm.dcd.main} \textbf{(c)} is proven.
\end{proof}

\begin{definition}
\label{def.perm.cycs.cycs}
\lean{AlgebraicCombinatorics.CycleDecomposition.cycles, AlgebraicCombinatorics.CycleDecomposition.cycleLengthsPartition}
\leantarget
\leanok
Let $X$ be a finite set. Let $\sigma$ be a
permutation of $X$. \medskip

\textbf{(a)} The \emph{cycles} of $\sigma$ are defined to be the cycles in the
DCD of $\sigma$ (as defined in Theorem \ref{thm.perm.dcd.main} \textbf{(a)}).
(This includes $1$-cycles, if there are any in the DCD of $\sigma$.)

We shall equate a cycle of $\sigma$ with any of its cyclic rotations; thus,
for example, $\left(  3,1,4\right)  $ and $\left(  1,4,3\right)  $ shall be
regarded as being the same cycle (but $\left(  3,1,4\right)  $ and $\left(
3,4,1\right)  $ shall not). \medskip

\textbf{(b)} The \emph{cycle lengths partition} of $\sigma$ shall denote the
partition of $\left\vert X\right\vert $ obtained by writing down the lengths
of the cycles of $\sigma$ in weakly decreasing order.
\end{definition}

\subsection{Properties of the cycle decomposition}

\begin{lemma}
\label{lem.cycles_pairwise_disjoint}
\lean{AlgebraicCombinatorics.CycleDecomposition.cycles_pairwise_disjoint}
\leanhelper
\leanok
The cycles of $\sigma$ are pairwise disjoint, i.e., their supports are pairwise disjoint.
\end{lemma}

\begin{proof}
\leanok
This is the cycle factors pairwise disjointness lemma from Mathlib.
\end{proof}

\begin{lemma}
\label{lem.cycles_noncommProd_eq}
\lean{AlgebraicCombinatorics.CycleDecomposition.cycles_noncommProd_eq}
\leanhelper
\leanok
The product of the cycles of $\sigma$ (in any order, since they commute) equals $\sigma$.
\end{lemma}

\begin{proof}
\leanok
The cycles are pairwise disjoint, hence pairwise commuting.
Their product equals $\sigma$ by the cycle factorization theorem from Mathlib.
\end{proof}

\begin{lemma}
\label{lem.cycleOf_mem_cycles}
\lean{AlgebraicCombinatorics.CycleDecomposition.cycleOf_mem_cycles}
\leanhelper
\leanok
If $x$ is not a fixed point of $\sigma$ (i.e., $\sigma(x) \neq x$),
then the cycle of $\sigma$ containing $x$ is among the cycles of $\sigma$.
\end{lemma}

\begin{proof}
\leanok
This is the characterization of cycle factor membership from Mathlib.
\end{proof}

\begin{lemma}
\label{lem.cycleType_eq_iff_isConj}
\lean{AlgebraicCombinatorics.CycleDecomposition.cycleType_eq_iff_isConj}
\leanhelper
\leanok
Two permutations $\sigma, \tau$ of a finite set $X$ have the same cycle type
(i.e., the same multiset of cycle lengths) if and only if they are conjugate, i.e.,
there exists a permutation $g$ of $X$ such that $\tau = g \sigma g^{-1}$.
\end{lemma}

\begin{proof}
\leanok
This is the conjugacy-iff-same-cycle-type characterization from Mathlib.
\end{proof}

\begin{definition}
\label{def.cycleLengths}
\lean{AlgebraicCombinatorics.CycleDecomposition.cycleLengths}
\leanhelper
\leanok
The \emph{cycle lengths} of $\sigma$ is the multiset of lengths
of the nontrivial cycles (length $\geq 2$) in the cycle decomposition of $\sigma$.
\end{definition}

\begin{lemma}
\label{lem.cycleLengths_sum}
\lean{AlgebraicCombinatorics.CycleDecomposition.cycleLengths_sum_eq_support_card}
\leanhelper
\leanok
The sum of the cycle lengths (of nontrivial cycles) equals the number of
non-fixed-point elements of $\sigma$.
\end{lemma}

\begin{proof}
\leanok
Follows from the cycle type summation lemma in Mathlib.
\end{proof}

\begin{lemma}
\label{lem.cycleLengths_mem_ge_two}
\lean{AlgebraicCombinatorics.CycleDecomposition.cycleLengths_mem_ge_two}
\leanhelper
\leanok
Every element of the cycle lengths multiset is at least $2$.
\end{lemma}

\begin{proof}
\leanok
Follows from the fact that each cycle type entry is at least $2$, which is proved in Mathlib.
\end{proof}

\begin{lemma}
\label{lem.cycleLengthsPartition_sum}
\lean{AlgebraicCombinatorics.CycleDecomposition.cycleLengthsPartition_sum}
\leanhelper
\leanok
The parts of the cycle lengths partition sum to $|X|$.
\end{lemma}

\begin{proof}
\leanok
This is the defining property of a partition of $|X|$.
\end{proof}

\subsection{Auxiliary lemmas on cycles}

\begin{lemma}
\label{lem.cyc_rotate_eq}
\lean{AlgebraicCombinatorics.CycleDecomposition.cyc_rotate_eq}
\leanhelper
\leanok
Let $l$ be a list with no duplicate entries, and let $k \in \mathbb{N}$.
Then the cycle permutation of $l$ equals the cycle permutation of the
$k$-fold rotation of $l$.
\end{lemma}

\begin{proof}
\leanok
This is the rotation invariance of cycle permutations from Mathlib.
\end{proof}

\begin{lemma}
\label{lem.cyc_eq_of_isRotated}
\lean{AlgebraicCombinatorics.CycleDecomposition.cyc_eq_of_isRotated}
\leanhelper
\leanok
Let $l_1, l_2$ be lists such that $l_1$ is a cyclic rotation of $l_2$,
and $l_1$ has no duplicate entries. Then $l_1$ and $l_2$ give rise to
the same cycle permutation.
\end{lemma}

\begin{proof}
\leanok
Follows from Lemma \ref{lem.cyc_rotate_eq}.
\end{proof}

\begin{proposition}
\label{prop.perm.cycs.same}
\lean{AlgebraicCombinatorics.CycleDecomposition.sameCycle_iff_exists_pow}
\leantarget
\leanok
Let $X$ be a finite set. Let $i,j\in X$. Let
$\sigma$ be a permutation of $X$. Then, we have the following chain of
equivalences:%
\begin{align*}
\  &  \left(  i\text{ and }j\text{ belong to the same cycle of }\sigma\right)
\\
\Longleftrightarrow\  &  \left(  i=\sigma^{p}\left(  j\right)  \text{ for some
}p\in\mathbb{N}\right) \\
\Longleftrightarrow\  &  \left(  j=\sigma^{p}\left(  i\right)  \text{ for some
}p\in\mathbb{N}\right)  .
\end{align*}
\end{proposition}

\begin{proof}
\leanok
The forward direction uses that for finite types, $\mathbb{Z}$-powers
can be replaced by $\mathbb{N}$-powers (reducing $k$ modulo the order of $\sigma$).
The backward direction is immediate.

The symmetric version (the third equivalence)
follows by applying the first equivalence to $\sigma^{-1}$.
\end{proof}

\begin{lemma}
\label{lem.mem_same_cycle_iff}
\lean{AlgebraicCombinatorics.CycleDecomposition.mem_same_cycle_iff}
\leanhelper
\leanok
Let $X$ be a finite set, $\sigma$ a permutation of $X$, and $i, j \in X$
with $i$ not a fixed point of $\sigma$. Then $i$ and $j$ belong to the
support of a common cycle of $\sigma$ if and only if $i$ and $j$ belong to the
same cycle of $\sigma$.
\end{lemma}

\begin{proof}
\leanok
For the forward direction, if $i$ and $j$ are both in the support of a cycle
$c$, then $i$ and $j$ are in the same cycle of $c$, and since $c$ is the cycle
of $\sigma$ containing $i$, this transfers to $i$ and $j$ being in the same
cycle of $\sigma$.
For the backward direction, the cycle of $\sigma$ containing $i$ is among
the cycles of $\sigma$, and both $i$ and $j$ belong to its support.
\end{proof}

\subsection{Sign from cycle decomposition}

\begin{lemma}
\label{lem.numCyclesTotal}
\lean{AlgebraicCombinatorics.CycleDecomposition.numCyclesTotal}
\leanhelper
\leanok
The total number of cycles of a permutation $\sigma$ of a finite set $X$
(including $1$-cycles / fixed points) equals the number of nontrivial cycles
(those of length $\geq 2$) plus the number of fixed points of $\sigma$,
i.e., the number of elements $x \in X$ with $\sigma(x) = x$.
\end{lemma}

\begin{proof}
\leanok
Immediate from the definitions.
\end{proof}

\begin{lemma}
\label{lem.sign_cycle}
\lean{AlgebraicCombinatorics.CycleDecomposition.sign_cycle}
\leanhelper
\leanok
Let $\sigma \in S_n$ be a $k$-cycle (i.e., $\sigma$ is a cycle and moves
exactly $k$ elements). Then
$(-1)^{\sigma} = (-1)^{k-1}$.
\end{lemma}

\begin{proof}
\leanok
This follows from the sign-of-a-cycle lemma in Mathlib, which gives
$(-1)^{\sigma} = -(-1)^{k}$ for a $k$-cycle $\sigma$.
Since $-(-1)^k = (-1)^{k-1}$ for $k \geq 1$, the result follows.
\end{proof}

\begin{proposition}
\label{prop.perm.cycs.sign}
\lean{AlgebraicCombinatorics.CycleDecomposition.sign_eq_neg_one_pow_card_sub_numCycles}
\leantarget
\leanok
Let $n\in\mathbb{N}$. Let $\sigma\in S_{n}$. Let
$k\in\mathbb{N}$ be such that $\sigma$ has exactly $k$ cycles (including the
$1$-cycles). Then, $\left(  -1\right)  ^{\sigma}=\left(  -1\right)  ^{n-k}$.
\end{proposition}

\begin{proof}
\leanok
Let
\begin{align*}
&  \left(  a_{1,1},a_{1,2},\ldots,a_{1,n_{1}}\right)  ,\\
&  \left(  a_{2,1},a_{2,2},\ldots,a_{2,n_{2}}\right)  ,\\
&  \ldots,\\
&  \left(  a_{k,1},a_{k,2},\ldots,a_{k,n_{k}}\right)
\end{align*}
be the $k$ cycles of $\sigma$. Thus,
\begin{align*}
\sigma &  =\operatorname*{cyc}\nolimits_{a_{1,1},a_{1,2},\ldots,a_{1,n_{1}}%
}\circ\operatorname*{cyc}\nolimits_{a_{2,1},a_{2,2},\ldots,a_{2,n_{2}}}%
\circ\cdots\circ\operatorname*{cyc}\nolimits_{a_{k,1},a_{k,2},\ldots
,a_{k,n_{k}}}
\end{align*}
so that, by Proposition
\ref{prop.perm.sign.props} \textbf{(e)} and \textbf{(c)},
\begin{align*}
\left(  -1\right)  ^{\sigma}
&  =\left(  -1\right)  ^{n_{1}-1}\cdot\left(  -1\right)  ^{n_{2}-1}\cdot
\cdots\cdot\left(  -1\right)  ^{n_{k}-1}\\
&  =\left(  -1\right)  ^{\left(  n_{1}-1\right)  +\left(  n_{2}-1\right)
+\cdots+\left(  n_{k}-1\right)  }.
\end{align*}
Since each element of $\left[  n\right]  $ appears exactly once in the
DCD, we have $n_{1}+n_{2}+\cdots+n_{k}=n$. Hence,%
\[
\left(  n_{1}-1\right)  +\left(  n_{2}-1\right)  +\cdots+\left(
n_{k}-1\right)  =\underbrace{\left(  n_{1}+n_{2}+\cdots+n_{k}\right)  }%
_{=n}-\,k=n-k.
\]
Thus, $\left(  -1\right)
^{\sigma}=\left(  -1\right)  ^{n-k}$. This proves Proposition
\ref{prop.perm.cycs.sign}.
\end{proof}

\begin{lemma}
\label{lem.sign_eq_neg_one_pow_sum_cycle_lengths_minus_one}
\lean{AlgebraicCombinatorics.CycleDecomposition.sign_eq_neg_one_pow_sum_cycle_lengths_minus_one}
\leanhelper
\leanok
Let $\sigma \in S_n$. Let $m_1, m_2, \ldots, m_j$ be the lengths of the
nontrivial cycles (of length $\geq 2$) of $\sigma$. Then
$(-1)^{\sigma} = (-1)^{(m_1 - 1) + (m_2 - 1) + \cdots + (m_j - 1)}$.
\end{lemma}

\begin{proof}
\leanok
This follows from the sign-from-cycle-type formula in Mathlib by observing that
$(m_1 + \cdots + m_j) + j$ and $(m_1 + \cdots + m_j) - j$ have the
same parity (they differ by $2j$).
\end{proof}

\subsection{Canonical DCD construction helpers}

\begin{definition}
\label{def.cycleMinElem}
\lean{AlgebraicCombinatorics.CycleDecomposition.cycleMinElem}
\leanhelper
\leanok
Let $c$ be a cycle permutation (i.e., $c$ is a cycle in the sense of
Definition~\ref{def.perm.cycs.cycs}).
The \emph{minimum element} of $c$ is the smallest element in the support of $c$.
\end{definition}

\begin{definition}
\label{def.cycleToCanonicalList}
\lean{AlgebraicCombinatorics.CycleDecomposition.cycleToCanonicalList}
\leanhelper
\leanok
Let $c$ be a cycle permutation. The \emph{canonical list representation}
of $c$ is the list $(a, c(a), c^2(a), \ldots)$ where $a$ is the minimum
element of $c$ (Definition~\ref{def.cycleMinElem}).
\end{definition}

\begin{lemma}
\label{lem.cycleToCanonicalList_head_le}
\lean{AlgebraicCombinatorics.CycleDecomposition.cycleToCanonicalList_head_le}
\leanhelper
\leanok
The first element of the canonical list is minimal among all elements of
the list.
\end{lemma}

\begin{proof}
\leanok
The first element is the minimum element of $c$, and
every element in the canonical list belongs to the support of $c$, so
the minimum is $\leq$ every element.
\end{proof}

\begin{lemma}
\label{lem.cycleToCanonicalList_formPerm}
\lean{AlgebraicCombinatorics.CycleDecomposition.cycleToCanonicalList_formPerm}
\leanhelper
\leanok
The cycle permutation determined by the canonical list equals $c$.
\end{lemma}

\begin{proof}
\leanok
The canonical list is the orbit of the minimum element under $c$.
The cycle permutation of this list recovers $c$ because the minimum element
is moved by $c$.
\end{proof}

\begin{lemma}
\label{lem.cycleMinElem_mem_support}
\lean{AlgebraicCombinatorics.CycleDecomposition.cycleMinElem_mem_support}
\leanhelper
\leanok
The minimum element of a cycle $c$ belongs to the support of $c$.
\end{lemma}

\begin{proof}
\leanok
By definition, the minimum element of $c$ is the minimum of the support,
which is an element of the support.
\end{proof}

\begin{lemma}
\label{lem.cycleToCanonicalList_nodup}
\lean{AlgebraicCombinatorics.CycleDecomposition.cycleToCanonicalList_nodup}
\leanhelper
\leanok
The canonical list representation of a cycle has no duplicate entries.
\end{lemma}

\begin{proof}
\leanok
Follows from the no-duplicates property of permutation cycle lists in Mathlib.
\end{proof}

\begin{lemma}
\label{lem.mem_cycleToCanonicalList_iff}
\lean{AlgebraicCombinatorics.CycleDecomposition.mem_cycleToCanonicalList_iff}
\leanhelper
\leanok
An element $x$ is in the canonical list of $c$ if and only if $x$ is in the support of $c$.
\end{lemma}

\begin{proof}
\leanok
The canonical list enumerates the orbit of the minimum element under $c$,
which is exactly the support of $c$ (since $c$ is a cycle).
\end{proof}

\begin{lemma}
\label{lem.cycleToCanonicalList_disjoint}
\lean{AlgebraicCombinatorics.CycleDecomposition.cycleToCanonicalList_disjoint}
\leanhelper
\leanok
If two cycles $c$ and $d$ are disjoint (as permutations), then their canonical
lists are disjoint (as lists).
\end{lemma}

\begin{proof}
\leanok
By Lemma~\ref{lem.mem_cycleToCanonicalList_iff}, membership in the canonical list
is equivalent to membership in the support.  Since $c$ and $d$ are disjoint,
their supports are disjoint.
\end{proof}

\begin{definition}
\label{def.dcdListToPerm}
\lean{AlgebraicCombinatorics.CycleDecomposition.dcdListToPerm}
\leanhelper
\leanok
Given a list of lists (each representing a cycle), the corresponding
permutation is obtained by composing the cycle permutations
$\operatorname{cyc}_{a_1, a_2, \ldots, a_m}$ for each list
$(a_1, a_2, \ldots, a_m)$ in the input.
\end{definition}

\begin{lemma}
\label{lem.dcdListToPerm_cyclesList_eq}
\lean{AlgebraicCombinatorics.CycleDecomposition.dcdListToPerm_cyclesList_eq}
\leanhelper
\leanok
The product of the cycle permutations of the canonical list representations of
$\sigma$'s cycles equals $\sigma$.
\end{lemma}

\begin{proof}
\leanok
Each canonical list determines a cycle permutation that recovers the original cycle (by
Lemma~\ref{lem.cycleToCanonicalList_formPerm}), and the product of all cycles
equals $\sigma$ (by the cycle factorization theorem from Mathlib).
\end{proof}

\begin{lemma}
\label{lem.cyclesList_flatten_nodup}
\lean{AlgebraicCombinatorics.CycleDecomposition.cyclesList_flatten_nodup}
\leanhelper
\leanok
The flattened list of all canonical cycle lists has no duplicate entries.
\end{lemma}

\begin{proof}
\leanok
Each individual canonical list is nodup (Lemma~\ref{lem.cycleToCanonicalList_nodup}),
and different cycles have disjoint canonical lists
(Lemma~\ref{lem.cycleToCanonicalList_disjoint}).
\end{proof}

\begin{lemma}
\label{lem.cyclesList_flatten_eq_support}
\lean{AlgebraicCombinatorics.CycleDecomposition.cyclesList_flatten_eq_support}
\leanhelper
\leanok
An element $y$ appears in the flattened list of canonical cycle lists if and only if
$y$ is not a fixed point of $\sigma$.
\end{lemma}

\begin{proof}
\leanok
By Lemma~\ref{lem.mem_cycleToCanonicalList_iff}, each canonical list covers
exactly the support of its cycle, and the union of all cycle supports equals
$\sigma$'s support.
\end{proof}
